% ===== PRÉAMBULE CANONIQUE — à coller VERBATIM en tête de CHAQUE .tex =====
\documentclass[11pt,a4paper]{article}
\usepackage[utf8]{inputenc}\usepackage[T1]{fontenc}\usepackage{lmodern}
\usepackage[french]{babel}
\usepackage{amsmath,amssymb,mathtools}\usepackage{icomma}
\usepackage[version=4]{mhchem}          % \ce{...} : équations & formules
\usepackage{siunitx}\sisetup{locale=FR} % \qty{}{}, \si{}, \ang{}
\usepackage{chemfig}                    % \chemfig{...} : structures développées
\usepackage{xcolor}\usepackage{colortbl}
\usepackage{tikz}\usepackage{pgfplots}\pgfplotsset{compat=1.18}
\usepackage[most]{tcolorbox}\usepackage{enumitem}\usepackage{array,booktabs}
\usepackage[a4paper,margin=2.2cm]{geometry}
\usetikzlibrary{arrows.meta,calc,angles,quotes,positioning,patterns,backgrounds,shapes.geometric}
\definecolor{primary}{RGB}{222,49,99}\definecolor{primarydark}{RGB}{150,22,55}
\definecolor{defcol}{RGB}{0,114,178}\definecolor{methcol}{RGB}{52,92,148}
\definecolor{excol}{RGB}{0,135,90}\definecolor{attncol}{RGB}{200,40,55}
\tcbset{boxbase/.style={enhanced,breakable,boxrule=0.9pt,arc=2.5pt,
  left=3mm,right=3mm,top=2.3mm,bottom=2.3mm,fonttitle=\bfseries,coltitle=white}}
% --- boîtes TUTORIEL (noms EXACTS, ne pas renommer) ---
\newtcolorbox{cle}[1][]{boxbase,colback=defcol!6,colframe=defcol,
  title={\textsc{À retenir}\ifx\relax#1\relax\relax\else\textmd{ : \itshape #1}\fi}}
\newtcolorbox{methode}[1][]{boxbase,colback=methcol!7,colframe=methcol,sharp corners,
  title={\textsc{Méthode}\ifx\relax#1\relax\relax\else\textmd{ --- \itshape #1}\fi}}
\newtcolorbox{exemplebox}[1][]{boxbase,colback=excol!5,colframe=excol,
  title={\textsc{Exemple}\ifx\relax#1\relax\relax\else\textmd{ : \itshape #1}\fi}}
\newtcolorbox{piege}[1][]{boxbase,colback=attncol!6,colframe=attncol,
  title={\textsc{Piège}\ifx\relax#1\relax\relax\else\textmd{ : \itshape #1}\fi}}
\newtcolorbox{remarque}[1][]{boxbase,colback=black!4,colframe=black!45,
  title={\textsc{Remarque}\ifx\relax#1\relax\relax\else\textmd{ : \itshape #1}\fi}}
% --- boîtes EXERCICE / CORRIGÉ (noms EXACTS, auto-comptées) ---
\newtcolorbox[auto counter]{exo}[1][]{boxbase,colback=primary!4,colframe=primary,
  title={\textsc{Exercice}~\thetcbcounter\ifx\relax#1\relax\relax\else\textmd{ --- \itshape #1}\fi}}
\newtcolorbox[auto counter]{corrige}[1][]{boxbase,colback=excol!4,colframe=excol,
  title={\textsc{Corrigé}~\thetcbcounter\ifx\relax#1\relax\relax\else\textmd{ --- \itshape #1}\fi}}
\newcommand{\R}{\mathbb{R}}\newcommand{\N}{\mathbb{N}}\newcommand{\Z}{\mathbb{Z}}\newcommand{\ds}{\displaystyle}
% (un \usepackage{fancyhdr}+en-tête est possible ; il est ignoré côté web)
% =========================================================================
\begin{document}

\begin{exo}[la loi combinée $pV/T=\text{cste}$]
Un gaz occupe $V_1=\qty{0.040}{\meter\cubed}$ sous $p_1=\qty{2.0e5}{\pascal}$ à $T_1=\qty{200}{\kelvin}$.
On le comprime jusqu'à $V_2=\qty{0.010}{\meter\cubed}$ tout en portant sa température à
$T_2=\qty{400}{\kelvin}$.
\begin{enumerate}[label=\alph*)]
  \item Écris la relation reliant les deux états lorsque $p$, $V$ et $T$ varient tous (à $n$ constant).
  \item Isole littéralement la pression finale $p_2$.
  \item Calcule $p_2$.
\end{enumerate}
\end{exo}

\begin{exo}[transformation en deux étapes]
Un cylindre à piston contient un gaz dans l'état initial $V_0=\qty{6.0}{\litre}$, $p_0=\qty{1.0}{atm}$,
$T_0=\qty{300}{\kelvin}$. On réalise deux étapes successives :
\begin{enumerate}[label=\alph*)]
  \item \textbf{Étape 1 — isotherme} : on comprime le gaz jusqu'à $p_1=\qty{3.0}{atm}$ (température
        inchangée). Calcule le volume $V_1$.
  \item \textbf{Étape 2 — isobare} : à $p_1=\qty{3.0}{atm}$, on chauffe de \qty{300}{\kelvin} à
        \qty{450}{\kelvin}. Calcule le volume final $V_2$.
  \item Retrouve directement $V_2$ à partir de l'état initial à l'aide de la loi combinée.
\end{enumerate}
\end{exo}

\begin{exo}[dilatation isobare d'un gaz chauffé]
À pression constante, on chauffe une quantité fixe de gaz de $\theta_1=\qty{127}{\degreeCelsius}$ à
$\theta_2=\qty{527}{\degreeCelsius}$.
\begin{enumerate}[label=\alph*)]
  \item Convertis les deux températures en kelvins.
  \item Par quel facteur le volume est-il multiplié ?
  \item Si le volume initial est $V_1=\qty{5.0}{\litre}$, calcule le volume final $V_2$.
\end{enumerate}
\end{exo}

\begin{exo}[sécurité d'une bombe aérosol (isochore)]
Une bombe aérosol rigide contient un gaz sous $p_1=\qty{3.0}{bar}$ à $\theta_1=\qty{27}{\degreeCelsius}$.
Le volume ne change pas.
\begin{enumerate}[label=\alph*)]
  \item On la laisse au soleil : sa température atteint \qty{127}{\degreeCelsius}. Calcule la nouvelle
        pression $p_2$.
  \item À quelle température (en \si{\degreeCelsius}) la pression atteindrait-elle \qty{6.0}{bar} ?
  \item Explique en une phrase pourquoi l'étiquette interdit de la jeter au feu.
\end{enumerate}
\end{exo}

\end{document}
